% =========================
%         PAQUETES
% =========================
\usepackage[utf8]{inputenc}
\usepackage[T1]{fontenc}
\usepackage[spanish,es-nodecimaldot]{babel}
% Evita conflictos de los caracteres < y > activos de babel/spanish con TikZ, también al compilar capítulos individuales.
\AtBeginDocument{\shorthandoff{<>}}
\usepackage{geometry}
\geometry{left=2.05cm,right=2.05cm,top=2.65cm,bottom=2.75cm,headsep=0.55cm}

% Matemática
\usepackage{amsmath,amssymb,amsfonts,amsthm,mathtools}
\usepackage{mathrsfs}
\usepackage{bm}
\usepackage{physics}
\usepackage{braket}

% Gráficos, tablas y diagramas
\usepackage{graphicx}
\usepackage{xcolor}
\usepackage[nopatch=footnote]{microtype}
\usepackage{tikz}
\usepackage{enumitem}
\usepackage{array}
\usepackage{booktabs}
\usepackage{float}
\usepackage[font=small,labelfont=bf,labelsep=period]{caption}
\usepackage[most]{tcolorbox}

% Enlaces y estructura
\usepackage{hyperref}
\usepackage{subfiles}
\usepackage{fancyhdr}
\usepackage{titlesec}
\usepackage{setspace}
\usepackage{csquotes}
\usepackage{tensor}
\usepackage{etoolbox}
\usetikzlibrary{arrows.meta,positioning,shapes.geometric,calc,decorations.pathmorphing,backgrounds,fit}

% =========================
%       PALETA VISUAL FRÍA
% =========================
\definecolor{UdeCAzul}{HTML}{243447}
\definecolor{UdeCAzulClaro}{HTML}{F4F6F8}
\definecolor{UdeCDorado}{HTML}{5F6F7A}
\definecolor{UdeCGris}{HTML}{6B7280}
\definecolor{QIRojo}{HTML}{4B5563}
\definecolor{QIVerde}{HTML}{52616B}
\graphicspath{{./}{../}}

% =========================
%     CONFIGURACIÓN GENERAL
% =========================
\setstretch{1.065}
\setlength{\headheight}{15pt}
\setlength{\parskip}{0.34em}
\setlength{\parindent}{0pt}
\setlength{\abovedisplayskip}{0.62em}
\setlength{\belowdisplayskip}{0.62em}
\setlength{\abovedisplayshortskip}{0.42em}
\setlength{\belowdisplayshortskip}{0.42em}
\emergencystretch=2.5em
\allowdisplaybreaks[2]

\hypersetup{
    colorlinks=true,
    hypertexnames=false,
    linkcolor=UdeCAzul,
    citecolor=UdeCAzul,
    urlcolor=UdeCAzul,
    pdftitle={Tópicos de Teoría Cuántica de la Información: Apuntes de apoyo al curso},
    pdfauthor={José Rosas Sepúlveda}
}

% Listas sobrias
\setlist[itemize,1]{label={\small$\triangleright$},leftmargin=1.65em,itemsep=0.10em,topsep=0.20em}
\setlist[itemize,2]{label={\small$\circ$},leftmargin=1.6em,itemsep=0.04em,topsep=0.12em}
\setlist[enumerate,1]{label={\textbf{\arabic*.}},leftmargin=1.7em,itemsep=0.12em,topsep=0.22em}

% Numeración y títulos para clase book
\renewcommand\thechapter{\arabic{chapter}}
\renewcommand\thesection{\thechapter.\arabic{section}}
\renewcommand\thesubsection{\thesection.\arabic{subsection}}
\renewcommand\thesubsubsection{\thesubsection.\arabic{subsubsection}}
\renewcommand{\contentsname}{Índice}
\setcounter{tocdepth}{2}
% Ancho de las etiquetas numéricas en el índice para capítulos, secciones y subsecciones.
\makeatletter
\renewcommand*\l@chapter{\@dottedtocline{0}{0em}{2.3em}}
\renewcommand*\l@section{\@dottedtocline{1}{1.5em}{3.0em}}
\renewcommand*\l@subsection{\@dottedtocline{2}{3.8em}{4.0em}}
\makeatother

\titleformat{\chapter}
    {\Huge\bfseries\color{UdeCAzul}}
    {\thechapter.}{0.75em}{}
\titleformat{name=\chapter,numberless}
    {\Huge\bfseries\color{UdeCAzul}}
    {}{0em}{}
\titleformat{\section}
    {\Large\bfseries\color{UdeCAzul}}
    {\thesection.}{0.6em}{}
\titleformat{\subsection}
    {\large\bfseries\color{UdeCAzul}}
    {\thesubsection.}{0.6em}{}
\titleformat{\subsubsection}
    {\normalsize\bfseries\color{UdeCAzul}}
    {\thesubsubsection.}{0.6em}{}
\titlespacing*{\chapter}{0pt}{0.0ex plus 0.5ex minus 0.2ex}{1.5ex plus 0.25ex}
\titlespacing*{\section}{0pt}{2.0ex plus 0.5ex minus 0.2ex}{1.0ex plus 0.2ex}
\titlespacing*{\subsection}{0pt}{1.55ex plus 0.35ex minus 0.2ex}{0.65ex plus 0.15ex}
\titlespacing*{\subsubsection}{0pt}{1.1ex plus 0.25ex minus 0.15ex}{0.45ex plus 0.12ex}

% =========================
%        ENCABEZADO
% =========================
\pagestyle{fancy}
\fancyhf{}
\lhead{\small \SubtituloApunte}
\rhead{\small\color{UdeCGris}Quantum Information}
\cfoot{\small \thepage}
\renewcommand{\headrulewidth}{0.35pt}
\renewcommand{\footrulewidth}{0pt}

\fancypagestyle{plain}{
    \fancyhf{}
    \lhead{\small \SubtituloApunte}
    \rhead{\small\color{UdeCGris}Quantum Information}
    \cfoot{\small \thepage}
    \renewcommand{\headrulewidth}{0.35pt}
}

% =========================
%      ENTORNOS TEÓRICOS
% =========================
\theoremstyle{definition}
\newtheorem{definicion}{Definición}[chapter]
\newtheorem{ejemplo}[definicion]{Ejemplo}
\newtheorem{observacion}[definicion]{Observación}
\newtheorem{notacion}[definicion]{Notación}

\theoremstyle{plain}
\newtheorem{proposicion}[definicion]{Proposición}
\newtheorem{teorema}[definicion]{Teorema}
\newtheorem{lema}[definicion]{Lema}
\newtheorem{corolario}[definicion]{Corolario}

\theoremstyle{remark}
\newtheorem*{comentario}{Comentario}

% Entornos discretos: sin cajas de color; sólo una barra lateral tenue.
\tcolorboxenvironment{observacion}{
    enhanced,breakable,blanker,
    borderline west={0.7pt}{0pt}{UdeCGris!55},
    left=7pt,right=0pt,top=2pt,bottom=2pt,
    before skip=6pt,after skip=6pt
}
\tcolorboxenvironment{definicion}{
    enhanced,breakable,blanker,
    borderline west={0.7pt}{0pt}{UdeCAzul!50},
    left=7pt,right=0pt,top=2pt,bottom=2pt,
    before skip=6pt,after skip=6pt
}
\tcolorboxenvironment{ejemplo}{
    enhanced,breakable,blanker,
    borderline west={0.7pt}{0pt}{UdeCGris!50},
    left=7pt,right=0pt,top=2pt,bottom=2pt,
    before skip=6pt,after skip=6pt
}
\tcolorboxenvironment{notacion}{
    enhanced,breakable,blanker,
    borderline west={0.7pt}{0pt}{UdeCGris!50},
    left=7pt,right=0pt,top=2pt,bottom=2pt,
    before skip=6pt,after skip=6pt
}
\tcolorboxenvironment{proposicion}{
    enhanced,breakable,blanker,
    borderline west={0.7pt}{0pt}{UdeCAzul!50},
    left=7pt,right=0pt,top=2pt,bottom=2pt,
    before skip=6pt,after skip=6pt
}
\tcolorboxenvironment{teorema}{
    enhanced,breakable,blanker,
    borderline west={0.7pt}{0pt}{UdeCAzul!65},
    left=7pt,right=0pt,top=2pt,bottom=2pt,
    before skip=6pt,after skip=6pt
}

% =========================
%        COMANDOS ÚTILES
% =========================
\newcommand{\C}{\mathbb{C}}
\newcommand{\R}{\mathbb{R}}
\newcommand{\N}{\mathbb{N}}
\newcommand{\Hil}{\mathcal{H}}
\newcommand{\I}{\mathbb{I}}
\newcommand{\ii}{\mathrm{i}}
\newcommand{\ee}{\mathrm{e}}

\newcommand{\TrA}{\operatorname{Tr}}
\providecommand{\rank}{\operatorname{rank}}
\newcommand{\diag}{\operatorname{diag}}
\newcommand{\Span}{\operatorname{span}}
\newcommand{\supp}{\operatorname{supp}}

\newcommand{\kb}[2]{\ket{#1}\!\bra{#2}}
\newcommand{\proj}[1]{\ket{#1}\!\bra{#1}}
\providecommand{\op}[1]{\hat{#1}}

\newcommand{\trA}{\operatorname{Tr}_A}
\newcommand{\trB}{\operatorname{Tr}_B}
\newcommand{\trC}{\operatorname{Tr}_C}

\newcommand{\cnot}{\mathrm{CNOT}}
\newcommand{\SWAP}{\mathrm{SWAP}}
\newcommand{\CZ}{\mathrm{CZ}}

% =========================
%     DATOS DEL DOCUMENTO
% =========================
\newcommand{\TituloApunte}{Quantum Information}
\newcommand{\SubtituloApunte}{Tópicos de Teoría Cuántica de la Información}
\newcommand{\AutorApunte}{José Rosas Sepúlveda}
\newcommand{\FechaApunte}{Mayo de 2026}
